\documentclass[12pt]{article}
\usepackage{amsmath}
\usepackage{geometry}
\geometry{margin=1in}

\begin{document}

\section*{Particle Interactions with Strangeness 2}

\textbf{For each interaction:}

\begin{itemize}
    \item Use the information about the interaction type (strong or weak).
    \item Apply conservation of charge, baryon number and strangeness.
\end{itemize}

\vspace{1em}

\textbf{Q1}

\[
K^- + n \rightarrow \Lambda^0 + X
\]

\begin{itemize}
    \item $\Lambda^0$ is a baryon, strangeness = $-1$
    \item The interaction is strong
\end{itemize}

Determine the identity of particle $X$.

\vspace{3em}

\textbf{Q2}

\[
\pi^+ + n \rightarrow K^+ + X
\]

\begin{itemize}
    \item $\pi^+$ is a meson
    \item The interaction is strong
    \item $X$ is a baryon
\end{itemize}

Write out a possible quark structure for baryon $X$.

\vspace{3em}

\textbf{Q3}

\[
\Sigma^- \rightarrow n + X
\]

\begin{itemize}
    \item $\Sigma^-$ is a baryon, strangeness = $-1$
    \item The decay occurs via the weak interaction
\end{itemize}

Determine the identity of particle $X$.

\newpage

\textbf{Q4}

\[
\pi^- + p \rightarrow \Lambda^0 + X
\]

\begin{itemize}
    \item $\Lambda^0$ is a baryon, strangeness = $-1$
    \item The interaction is strong
    \item $X$ is a meson
\end{itemize}

Write out a possible quark structure for meson $X$.

\vspace{3em}

\textbf{Q5}

\[
K^0 + p \rightarrow \Sigma^+ + X
\]

\begin{itemize}
    \item $\Sigma^+$ is a baryon, strangeness = $-1$
    \item The interaction is strong
\end{itemize}

Determine the identity of particle $X$.

\vspace{3em}

\textbf{Q6}

\[
K^+ \rightarrow \pi^+ + X
\]

\begin{itemize}
    \item The decay occurs via the weak interaction
\end{itemize}

Determine the identity of particle $X$.

\vspace{3em}

\textbf{Q7}

\[
\Lambda^0 \rightarrow p + \mu^- + X
\]

\begin{itemize}
    \item $\Lambda^0$ is a baryon, strangeness = $-1$
    \item The decay occurs via the weak interaction
\end{itemize}

Determine the identity of particle $X$.

\newpage

\section*{Answer Key}

\begin{enumerate}
    \item $X = \pi^{-}$
    \item A possible quark structure for $X$ is $uds$
    \item $X = \pi^{-}$
    \item A possible quark structure for $X$ is $u\bar{u}$ or $d\bar{d}$ (e.g.\ $\pi^{0}$)
    \item $X = \pi^{+}$
    \item $X = \pi^{0}$
    \item $X = \bar{\nu}_\mu$
\end{enumerate}

\end{document}